\documentclass[11pt]{article}
% Standalone preamble for web-only supplementary notes.
\usepackage[T1]{fontenc}
\usepackage{amsmath,amssymb,amsthm}
\usepackage{hyperref}
\usepackage{cleveref}
\DeclareMathOperator{\Map}{Map}
\DeclareMathOperator{\Fun}{Fun}
\newcommand{\iso}{\xrightarrow{\sim}}
\newtheorem{lemma}{Supplementary Lemma}
\newtheorem{corollary}[lemma]{Supplementary Corollary}
\crefname{lemma}{Supplementary Lemma}{Supplementary Lemmas}
\Crefname{lemma}{Supplementary Lemma}{Supplementary Lemmas}
\crefname{corollary}{Supplementary Corollary}{Supplementary Corollaries}
\Crefname{corollary}{Supplementary Corollary}{Supplementary Corollaries}

\begin{document}
\section*{Lifting a family from three restrictions}

To recover the morphisms of a \href{subcategories.html}{generated subcategory},
the formalization uses the three endomorphisms of the walking morphism.
The calculation is about a whole family of arrows. Here is how its three
restrictions are assembled, using specified lifts rather than a choice
made separately at every point.

Assume the Chapter 1 mapping-anima and coproduct structures. For a functor
$m\colon X\to Y$, write $m_*$ for postcomposition on mapping animae.
A lift of $h$ through $m_*$ includes a functor $\widetilde h$ and an
identification $m_*\widetilde h=h$.

\begin{lemma}
\label[lemma]{lem:Supplementary_Summand_Lifting}
Let $h\colon\Gamma\to\Map(A\sqcup B,Y)$, where $\Gamma$ is any absolute
category. If the restrictions of $h$ to $A$ and $B$ have specified lifts
through $m_*$, then $h$ has a lift through
$\Map(A\sqcup B,X)\to\Map(A\sqcup B,Y)$.
\end{lemma}
\begin{proof}[Proof idea]
Pair the two supplied lifts and lift this pair through the equivalence
\[
 \Map(A\sqcup B,X)\iso\Map(A,X)\times\Map(B,X).
\]
Restriction to the two summands commutes with postcomposition by $m$.
Consequently, the two supplied comparison identifications give an
identification of the restrictions of $m_*\widetilde h$ and $h$.
Reflect it through the corresponding restriction equivalence for $Y$.
This gives the required comparison for the whole family. The argument
does not require $m$ to be an embedding.
\end{proof}

\begin{lemma}
\label[lemma]{lem:Supplementary_Three_Restrictions}
Suppose $x_0,x_1,x_2\colon\star\to K$ induce an equivalence
$(\star\sqcup\star)\sqcup\star\iso K$.
A family $h\colon\Gamma\to\Map(K,Y)$ lifts through $m_*$ if its
three restrictions to $\Map(\star,Y)$ have specified lifts.
\end{lemma}
\begin{proof}[Proof idea]
Apply \Cref{lem:Supplementary_Summand_Lifting} twice. At each step,
transport the supplied comparisons along the computation identifications
for the coproduct inclusions. This lifts the restriction of $h$ along
the displayed equivalence. Precomposition by that equivalence is itself
an equivalence on mapping animae. Lift back through it and reflect the
comparison with $h$. Notice that the three restrictions take values in
$\Map(\star,Y)=Y^\simeq$; we keep these core-valued restrictions explicit, rather than identifying
$\Map(\star,Y)$ with $Y$.
\end{proof}

\begin{lemma}
\label[lemma]{lem:Supplementary_Generated_Arrows}
Assume also the walking-morphism, endpoint, Segal, commutative-square,
interval-core and coproduct-universality structures of the preceding
chapters. Let $m\colon M\hookrightarrow\Map([1],C)$ be a collection
closed under source and target identities, and let $i\colon A\to C$
be a subcategory presentation of $M$. The comparison
\[
 \Map([1],A)\longrightarrow M
\]
is an equivalence.
\end{lemma}
\begin{proof}[Proof idea]
Use the established three-point presentation of $\Map([1],[1])$,
whose points name the constant functor at $0$, the identity functor,
and the constant functor at $1$.
The family $M\to\Map([1],C)$ acts by postcomposition on this anima.
Its three restrictions are the source identity, the given arrow, and the
target identity. The two closure witnesses supply the first and third
lifts; the middle lift is the identity. The naming and evaluation
comparisons put these lifts in the form needed by
\Cref{lem:Supplementary_Three_Restrictions}.

The resulting lift of the whole action supplies a cone for the defining
subcategory pullback, tested at $[1]$. Its factorization gives a section
of $\Map([1],A)\to M$. The comparison is an embedding because both maps
into $\Map([1],C)$ are embeddings. An embedding with a section is an
equivalence. Closure under composition is needed to obtain the presentation
from the subcategory axiom, but this calculation, with a presentation
already supplied, uses only closure under identities.
\end{proof}
\end{document}
